% page 567 du fac-simile (image p-566 du scan)
% bloc 185.4 x 259.3 mm ; marge gauche 16.6 mm
\pgc{-5.080mm}{0.000mm}{
\l{}
\l{}
\l{}
\l{}
\l{}
\l{}
\l{}
\l{\sou{submisionar}. (\sou{netrans., pri}).Ofrar sua vari, sua laboro, a la admi-}
\l{\cel{3}nistrantaro statala, municipala, e c., lor adjudiko-kompro. -\cel{2}DF.}
\l{}
\l{\sou{subornar}. (\sou{trans.}) Instigar ad ago kontre-etika, mala-mora, per lurivo,}
\l{\cel{3}per charmivo. - EFIS.}
\l{}
\l{\sou{subrek}argo. (\sou{komerco per navo-transporto)}. Prokuraciero di la persono}
\l{\cel{3}o personi qua o qui proprietas la vari embarkita sur navo. -}
\l{}
\l{\sou{subreptar}. (\sou{trans}.) (\sou{yurocienco)}. Obtenar per moyeni frauda, kontre-}
\l{\cel{3}lega e trompa. - EFIS.}
\l{}
\l{\sou{subreto}. (\sou{teatro}). En komedio, la rolo (acesora) di chambro-servistino,}
\l{\cel{3}di akompanantino, di qua la karaktero esas impetuoza, audacoza, tre}
\l{\cel{3}alerta, tre inteligenta. - DEFR.}
\l{}
\l{\sou{subrogar}. (\sou{trans.}) (\sou{yurocienco}). Establisar en la rolo di ulu altru qua}
\l{\cel{3}esis la tituliero normala. - DEFIS.}
\l{}
\l{\sou{subsidiar}. (\sou{trans., pro. pri)}. Donacar pekunio ad ulu qua indijas ex-}
\l{\cel{3}treme, e qua bezonas lu por exekutendi qui urjas. (Kom ex. Subsidiar}
\l{\cel{3}navigado-kompanio pri la konstrukto di paketboti nova-tipa). - DEFIRS.}
\l{}
\l{\sou{substanco}. I. To ye quo konsistas la ento. - II. (\sou{filoz}.)To quo per-}
\l{\cel{3}manas kontinue en ento. - III. Materio ye qua kozo konsistas.-DEFIS.}
\l{}
\l{\sou{substantivo}. (\sou{gram.}) Vorto qua indikas la ento (kontraste kun la\cel{1}\sou{adjek}-}
\l{\cel{3}tivo,qua indikas la eso-maniero).-\cel{2}DEFIS.}
\l{}
\l{\sou{substitucar}. (\sou{trans., ad.}) Pozar (persono, kozo) en la plaso di altra.}
\l{\cel{3}- DEFIS.}
\l{}
\l{\sou{substituto}. (\sou{yurocienco}). Judiciisto, en tribunalo, qua komisesis}
\l{\cel{3}por agar vice altra qua esas absenta od impedata. - DEFIS.}
\l{}
\l{\sou{subsumar}. (\sou{trans.}). (\sou{filoz.}) I.Situar logikale individuo o speco}
\l{\cel{3}aden la klaso o genero qua implikas li.- II.(\sou{che Kant}). Pensar}
\l{\cel{3}intuice sub koncepto o kategorio di intelekto.- DEFIS.}
\l{}
\l{subterfujo. Moyeno nedireta di eskapar ulo embarasiva, evitar}
\l{\cel{3}responso o justifikivo. - EFIS.}
\l{}
\l{\sou{subtila.}\cel{1}I. Qua konsistas ye elementi extreme tenua. - II. (\sou{metaf.})}
\l{\cel{3}Di qua la nuanci pensala esas tenuega, e dicernebla nur desfacile.}
\l{\cel{3}- III. Qua perceptas la nuanci pensala, mem tre desfacile dicer-}
\l{\cel{3}nebla. - DEFIS.}
\l{}
\l{subvencionar. (\sou{trans.)}\cel{2}Grantar (ad ulu, ad ulo) parto ek la pekunio-}
\l{\cel{3}stoko qua esas rezervita por spensi neexpektita, por helpar finance}
\l{\cel{4}entraprezajo. - DEFIS.}
\l{}
\l{subversar.\cel{2}(trans.)\cel{2}uaze renversar, en la spiriti, omna lego, omna}
\l{\cel{4}regulo, omna principo, generale agnoskata kom bona. - DEFIS.}
}
